\scanpage{79}
\mspage{75}
\deleted{\unsure{À|A}} Un spatios est \deleted{associé} \inserted{muni} d'une structure de disjonction canonique :
\[
x\mathrel{|\unsure{\circ|0}|}y\mathrel{\overset{\mathrm{def}}{\Longleftrightarrow}}x\leq Cy\quad(\Longleftrightarrow y\leq Cx,\ \text{par Spat 2}).
\]
\unsure{Notons|Noter|Observons} que
\[
x\mathrel{|\circ|}y,\ x'\leq x,\ y'\leq y\Longrightarrow x'\mathrel{|\circ|}y'.
\]
% slanted
\inserted{\uline{NB} Si $x=\operatorname{Sup}x_{i}$, $y=\operatorname{Sup}y_{j}$, $x_{i}\mathrel{|\circ|}y_{j}\quad\forall i,j$, alors $x\mathrel{|\circ|}y$ (et inversement)}

\deleted{La structure de disjonction permet de récupérer la structure $\leq,C$, grâce à la}

\deleted{\uline{Proposition} Soit $\Sigma$ un spatios, considérons}

\unsure{1|l'} \uline{Exemple de Spatios} Si $\underline{E}=(E,\mathrel{|\circ|})$ est un disjoncteur, $(\mathbb{\Sigma}_{\underline{E}},\subset,C)$ est un spatios. Les \deleted{\unsure{relations}} propriétés Spat 1, Spat 2 ont été dégagées, elles ne dépendent que du fait que $\mathrel{|\circ|}_{E}$ soit symétrique \inserted{(Disj 1)}, sans plus. La propriété Spat 3 résulte de (Disj. 2), qui \unsure{implique|signifie}, pour $\forall A\in\unsure{\mathcal{P}|\mathrm{P}}(E)$
\[
A\cap\operatorname{cosupp}_{E}^{\circ}(A)\subset E_{0}
\]
donc, si $A\supset E_{\unsure{0}}$, et en particulier si $A\unsure{\in|\subset}\mathbb{\Sigma}_{E}$
\[
A\cap\operatorname{cosupp}_{E}^{\circ}(A)=E_{0}.
\]

% slanted
\margin{\uline{Ensemble ordonné disjonctif} ou \uline{disjoncteur ordonné}. $(E,\overset{\circ}{\ll},\mathrel{|\circ|})$ avec $x\mathrel{|\circ|}y,\ x'\mathrel{\overset{\circ}{\ll}}x,\ y'\mathrel{\overset{\circ}{\ll}}y\Longrightarrow x'\mathrel{|\circ|}y'$. Ex : Tout spatios est un disjoncteur ordonné. Si $(E,\mathrel{\overset{\circ}{\ll}},\mathrel{|\circ|})$ est un disjoncteur ordonné, $\sigma_{E}:E\longrightarrow\mathbb{\Sigma}_{\underline{E}}$ est \uline{croissante}.}

% slanted
\margin{\uline{NB} Tout disj. \deleted{\unsure{ordonné}} est muni d'une relation d'ordre \uline{canonique}, qui en fait un disj. ordonné (la \unsure{moins} fine de toutes celles qui sont \unsure{comp.} avec la str. de disjonction) : \unsure{à savoir|c'est|à dire} l'image inverse de la structure d'ordre de $\mathbb{\Sigma}_{E}$.}

Ainsi, la structure de disjonction sur un \unsure{disjoncteur|ensemble} quelconque, \unsure{se|on} déduit de celle sur le \deleted{disjoncteur} spatios associé $\mathbb{\Sigma}_{\underline{E}}$, via l'application canonique
\[
\sigma_{\underline{E}}:E\longrightarrow\mathbb{\Sigma}_{\underline{E}}.
\]
% archive sheet number: 78
